% RESEARCH_DRAFT_SHA256: caac53fc4091447ad127eaa2b10f23b70231f7e99d7745422ad312b32ccc58ff
% Layout and source-reference wrappers expanded for the web renderer.
\documentclass[11pt]{article}
\usepackage[margin=25mm]{geometry}
\usepackage{amsthm}
\usepackage{fontspec,unicode-math}
\setmainfont{Latin Modern Roman}
\setmathfont{Latin Modern Math}
\usepackage{amsmath,amssymb,mathtools,xurl,hyperref}
\hypersetup{colorlinks=true,urlcolor=blue,linkcolor=blue}
\setcounter{secnumdepth}{0}
\setlength{\parindent}{0pt}
\setlength{\parskip}{5pt}
\setlength{\emergencystretch}{3em}
\newcommand{\sref}[2]{\hyperlink{source:#1:#2}{[S#2]}}
\newcommand{\cataloguescope}[1]{\paragraph{Source record.}#1}
\newcommand{\floor}[1]{\left\lfloor#1\right\rfloor}
\newtheorem{proposition}{Proposition}
\newtheorem{lemma}[proposition]{Lemma}
\begin{document}
\paragraph{Research manuscript.} OpenAI \texttt{gpt-6-astra}, reasoning \texttt{ultra}; 11 October 2026.
\paragraph{Preserved source body.} The complete catalogue mathematical text and sources follow with layout and source links expanded, including source qualifications and its historical preparation notes. New research begins at the section headed Research attempt.

\section{Observation sufficiency for non-computational mathematical work}\label{statement}

\subsection{Definitions and mathematical statement}

\paragraph{Shared notation.}Standard number systems are $\mathbb N,\mathbb Z,\mathbb Q,\mathbb R,\mathbb C$. Specialized notation and index ranges follow the cited source and the definitions below.


A computational exercise asks a learner to carry out a calculation or algorithm. Mathematical activities can also involve constructing definitions, proposing conjectures, explaining reasoning, finding counterexamples, or proving claims. An assessment task needs a stated learning objective and a way to interpret the work it elicits.
\paragraph{Record scope and qualifications.}
The catalogue title is retained as metadata. Named certificates, contracts, bounds or frameworks in that title are not asserted as results of the original source extract.
The phrase task oriented is not operationally defined, and the registry marks the wording as reconstructed. No replacement taxonomy or theorem is imposed.

Catalogue screening limitation: Educational assessment, adoption, or study-design discussion; the supplied question does not state a mathematical target with defined assumptions and outcome.

The source snapshot labels this record partially solved; this is catalogue metadata rather than independent verification of a mathematical result or the current completion of a project.
The source text below is rendered literally where necessary; damaged formulas and composite extracts are not silently repaired into a new mathematical claim.
The supplied catalogue marks this formulation as an unverified reconstruction; missing assumptions are not certified by the source link.
\paragraph{Statement recorded in the supplied source.}
35. Problems which are not "task oriented" (e.g. not computational) \href{http://aimath.org/pastworkshops/oerassessmentproblems.pdf}{[S2]}
\paragraph{Source status.}
The supplied catalogue labels this source record partially solved; that metadata was checked against the supplied snapshot on 2026-10-11. The label is preserved for provenance and is not an independent verification of present open status or mathematical correctness.
\subsection{Short English statement}
Explore mathematical activities beyond narrowly computational or task-oriented exercises.
\subsection{Sources}
\begin{itemize}
\item[S1] Ulam AI and the UnsolvedMath Contributors, UnsolvedMath: A Curated Collection of Open Mathematics Problems, catalogue entry. Catalogue curation and metadata: CC BY 4.0. The source extract is retained; definitions, English summaries and qualifications are editorial additions.
\url{https://huggingface.co/datasets/ulamai/UnsolvedMath}.
\item[S2] AIM Workshop Problem Lists, source linked by the catalogue.
\url{http://aimath.org/pastworkshops/oerassessmentproblems.pdf}.
\end{itemize}
\paragraph{Preparation and review.}
The original source and UnsolvedMath contributors retain attribution for the source extract. Codex (GPT-6; exact serving revision not exposed) assisted with definitions, English summaries, scope notes and formatting. No new human mathematical authorship or independently verified solution is asserted. Catalogue identity checked against the complete supplied registry and current repository. Distinct upstream records, including repeated questions and broad or damaged extracts, are retained for complete coverage of the source CS category. A repeated record is not a new independent conjecture. Complete mathematical-equivalence and current-literature review has not been performed. Preparation and final review follow contribution prompts C and D, with source gaps disclosed. Structural and rendering checks do not establish a new theorem or open status.
\label{end:statement}

\section{Research attempt}
\subsection{Literature, source context and scope}
The primary AIM notes describe a brainstorming session. The original excerpt and all catalogue qualifications remain above. They do not specify a mathematical theorem or a complete empirical study. The formal model below is an explicitly limited research route introduced for this attempt, not a reconstruction of a missing original conjecture. The supplied catalogue's earlier research summary was inspected as context; its named framework is not adopted as an established solution. The elementary result below is not claimed novel, and no educational effectiveness data were collected.

\subsection{Attempt: assessing a proof certificate rather than only a final answer}
The phrase non-task-oriented is not defined by the source, and its reconstruction remains unverified. I tested a restricted representation for non-computational work: a finite directed graph of claims, with edges from cited premises to conclusions, checked against declared axioms and sound, decidably checkable inference rules. Sound means that true premises imply a true conclusion in every intended interpretation. Here a leaf means a node with no incoming premise edges.
\begin{proposition}
If the graph is acyclic, every leaf is a verified axiom, and every nonleaf is a valid inference from its incoming premises, then every claim in the graph is true in each interpretation satisfying the axioms. Verification takes $O(|V|+|E|)$ graph work plus the costs of checking axioms and individual inference steps. Dropping acyclicity invalidates this guarantee.
\end{proposition}
\begin{proof}
Topologically order the finite acyclic graph. Such an order exists: a finite acyclic directed graph has a vertex with no incoming edge, since repeatedly following predecessors otherwise repeats a vertex and forms a cycle; remove such a vertex and continue. Processing each vertex and edge once with incoming-edge counts computes an order in linear time. In that order, a leaf is true by the axiom check. Every later claim has already established true premises, so soundness of its verified rule makes it true. Induction proves the result for all vertices. Local checking costs are additional and are not assumed uniformly small.

For failure without acyclicity, use a single node labeled $0=1$ with an edge to itself and the identity rule ``from $P$ infer $P$.'' That rule is sound and its local premise/conclusion match is valid, but the resulting circular object has no axiom foundation and its conclusion is false in ordinary arithmetic. Acyclicity rejects it.
\end{proof}
\subsection{Verification and remaining gaps}
A topological checker accepts a three-node axiom-to-inference chain and rejects the one-node self-loop. These finite checks illustrate the structural distinction. The result is the standard soundness induction for a proof certificate, not a universal natural-language grader. It does not evaluate mathematical insight, explanatory quality, authorship, understanding or all forms of creative work, and it depends on the chosen sound rule system. The original pedagogical agenda remains broader and undefined as a mathematical target.

\section{References and source attribution}
American Institute of Mathematics, workshop discussion notes on open educational resources and assessment, \url{https://aimath.org/pastworkshops/oerassessmentproblems.pdf}. The 16-page primary PDF was opened on 11 October 2026 and its numbered brainstorming item was checked against the supplied record. The source supports the provenance and agenda, not the additional modeling assumptions or an empirical educational conclusion. Catalogue and contributor attribution are retained in the reproduced Sources section. The worked mathematics is proved here without reliance on an unverified external theorem.


\paragraph{Integration review.} Codex (GPT-6; exact serving revision not exposed) reviewed the written arguments and expanded layout and source-reference wrappers. Repository metadata and the production web renderer were checked using Python 3.13 and JavaScript with MathJax 3.2.2. This model-assisted review is not independent human proof verification; no PDF compilation or formal proof checking is asserted.

\section{Completion Estimate}
\noindent\textbf{COMPLETION: 0\%}
No substantiated new progress toward the entire original target is claimed. The worked results and failed route are delimited above; a known result, illustrative model, or auxiliary obstruction is not counted as original-target completion.

\section{Final status and attribution}
\textbf{Attempt status: unresolved.} This describes this research attempt, not a new ruling on the source catalogue's status or on external literature claims.
The mathematics and this self-audit were generated by OpenAI \texttt{gpt-6-astra}, reasoning \texttt{ultra}, on 11 October 2026. Codex (GPT-6; exact serving revision not exposed) orchestrated the batch and contributed formatting/repository review. No human mathematical authorship, independent human proof checking, or formal verification is claimed. The historical preparation paragraphs in the reproduced catalogue body concern that body only.

Contribution prompt A was applied to the complete definition and supplied source record, followed by the adversarial checks of prompt B. Local reproduction files are \texttt{batch008/checks.py}, \texttt{computations.json}, and the per-ID audit JSON. They were not public links at preparation time. Numerical checks are finite and do not prove unrestricted results. No LaTeX engine was found locally; compilation was not performed by this research agent.

\end{document}
